\documentclass[11pt,oneside]{book}

\usepackage[a4paper,margin=1in]{geometry}
\usepackage{amsmath,amssymb,amsthm,mathrsfs}
\usepackage{graphicx}
\usepackage{hyperref}
\usepackage{physics}
\usepackage{bm}
\usepackage{enumitem}
\usepackage{tensor}
\usepackage{mathtools}

\hypersetup{
    colorlinks=true,
    linkcolor=blue,
    citecolor=blue,
    urlcolor=blue
}

\newtheorem{theorem}{Theorem}[chapter]
\newtheorem{proposition}[theorem]{Proposition}
\newtheorem{lemma}[theorem]{Lemma}
\newtheorem{definition}[theorem]{Definition}
\newtheorem{remark}[theorem]{Remark}

\title{Classical and Quantum Behavior of the Double Pendulum\\
\large Geometry, Chaos, and Semiclassical Analysis}

\author{}
\date{}

\begin{document}

\frontmatter
\maketitle
\tableofcontents

\mainmatter

%%%%%%%%%%%%%%%%%%%%%%%%%%%%%%%%%%%%%%%%%%%%%%%%%%%%%%%%%%%%
\chapter{Introduction}

The double pendulum is one of the simplest mechanical systems exhibiting
nonlinearity, non-integrability, and deterministic chaos.
Despite its apparent simplicity, it connects:

\begin{itemize}
\item Riemannian geometry,
\item Hamiltonian mechanics,
\item symplectic topology,
\item nonlinear dynamics,
\item spectral theory,
\item quantum chaos.
\end{itemize}

The goal of this monograph is to present a complete treatment
from classical foundations to semiclassical quantization.

%%%%%%%%%%%%%%%%%%%%%%%%%%%%%%%%%%%%%%%%%%%%%%%%%%%%%%%%%%%%
\chapter{Configuration Manifold and Kinematics}

\section{Physical Model}

Two point masses $m_1,m_2$ are connected by rigid rods
$\ell_1,\ell_2$ constrained to a plane.

Generalized coordinates:
\[
q=(\theta_1,\theta_2), \quad \theta_i\in S^1.
\]

\section{Configuration Space}

\[
Q=S^1\times S^1 = T^2.
\]

\section{Phase Space}

\[
M=T^*Q.
\]

Canonical coordinates:
\[
(\theta_1,\theta_2,p_1,p_2).
\]

Symplectic form:
\[
\omega=d\theta_1\wedge dp_1+d\theta_2\wedge dp_2.
\]

%%%%%%%%%%%%%%%%%%%%%%%%%%%%%%%%%%%%%%%%%%%%%%%%%%%%%%%%%%%%
\chapter{Lagrangian and Hamiltonian Structure}

\section{Kinetic Energy}

\[
T=\frac12 m_1 \ell_1^2 \dot\theta_1^2
+\frac12 m_2(\ell_1^2\dot\theta_1^2
+\ell_2^2\dot\theta_2^2
+2\ell_1\ell_2\dot\theta_1\dot\theta_2\cos\Delta),
\]
where $\Delta=\theta_1-\theta_2$.

\section{Potential}

\[
V=-(m_1+m_2)g\ell_1\cos\theta_1
-m_2 g\ell_2\cos\theta_2.
\]

\section{Metric Form}

Define:

\[
g_{ij}=
\begin{pmatrix}
A & B\cos\Delta\\
B\cos\Delta & C
\end{pmatrix}
\]

where
\[
A=(m_1+m_2)\ell_1^2,
\quad
B=m_2\ell_1\ell_2,
\quad
C=m_2\ell_2^2.
\]

Hamiltonian:
\[
H=\frac12 g^{ij}p_i p_j + V(\theta).
\]

%%%%%%%%%%%%%%%%%%%%%%%%%%%%%%%%%%%%%%%%%%%%%%%%%%%%%%%%%%%%
\chapter{Curvature of the Configuration Metric}

\section{Determinant}

\[
|g|=AC-B^2\cos^2\Delta.
\]

\section{Gaussian Curvature}

\begin{theorem}
The Gaussian curvature of the configuration metric is
\[
K(\Delta)
=
\frac{B^2 \sin^2\Delta
(AC+B^2\cos^2\Delta)}
{2(AC-B^2\cos^2\Delta)^2}.
\]
\end{theorem}

\begin{proof}
Compute Christoffel symbols via
\[
\Gamma^k_{ij}=\frac12 g^{kl}
(\partial_i g_{jl}
+\partial_j g_{il}
-\partial_l g_{ij}),
\]
observe dependence only on $\Delta$,
and substitute into standard 2D curvature formula.
\end{proof}

\begin{remark}
Curvature vanishes at $\Delta=0,\pi$ and is maximal at $\Delta=\pi/2$.
\end{remark}

%%%%%%%%%%%%%%%%%%%%%%%%%%%%%%%%%%%%%%%%%%%%%%%%%%%%%%%%%%%%
\chapter{Linear and Nonlinear Dynamics}

\section{Small Oscillations}

Linearization yields matrix equation:

\[
M\ddot{\Theta}+K\Theta=0.
\]

Normal mode frequencies follow from:

\[
\det(K-\omega^2 M)=0.
\]

\section{Energy Surfaces}

Energy hypersurfaces:

\[
H(\theta,p)=E.
\]

Topology changes with $E$.

%%%%%%%%%%%%%%%%%%%%%%%%%%%%%%%%%%%%%%%%%%%%%%%%%%%%%%%%%%%%
\chapter{Chaos and Lyapunov Exponents}

\section{Tangent Flow}

\[
\dot{\delta z}=DX_H(z)\delta z.
\]

Largest Lyapunov exponent:

\[
\lambda=\lim_{t\to\infty}
\frac1t\ln\frac{\|\delta z(t)\|}{\|\delta z(0)\|}.
\]

\section{Poincar\'e Sections}

Fix $\theta_2=0$, $p_2>0$.
Resulting 2D map reveals invariant curves and chaotic regions.

%%%%%%%%%%%%%%%%%%%%%%%%%%%%%%%%%%%%%%%%%%%%%%%%%%%%%%%%%%%%
\chapter{Symplectic Numerical Integration}

\section{Implicit Midpoint Rule}

\[
\frac{z_{n+1}-z_n}{h}
=
X_H\left(\frac{z_{n+1}+z_n}{2}\right).
\]

Second-order, symplectic, time-reversible.

\section{Composition to Fourth Order}

\[
\Phi_h^{(4)}
=
\Phi_{\alpha h}^{(2)}
\circ
\Phi_{\beta h}^{(2)}
\circ
\Phi_{\alpha h}^{(2)},
\]

\[
\alpha=\frac1{2-2^{1/3}}, \quad
\beta=1-2\alpha.
\]

\section{Lyapunov Algorithm}

Renormalize tangent vector periodically:
\[
\lambda\approx \frac1{T}
\sum_k \ln \frac{\|\delta z_k\|}{\|\delta z_{k-1}\|}.
\]

%%%%%%%%%%%%%%%%%%%%%%%%%%%%%%%%%%%%%%%%%%%%%%%%%%%%%%%%%%%%
\chapter{Quantum Double Pendulum}

\section{Hilbert Space}

\[
\mathcal H=L^2(T^2).
\]

\section{Hamiltonian Operator}

\[
\hat H=-\frac{\hbar^2}{2}\Delta_g+V.
\]

\begin{theorem}
On compact $T^2$, $\hat H$ is essentially self-adjoint.
\end{theorem}

%%%%%%%%%%%%%%%%%%%%%%%%%%%%%%%%%%%%%%%%%%%%%%%%%%%%%%%%%%%%
\chapter{Semiclassical Analysis}

\section{Van Vleck Propagator}

\[
K\sim
\sum_{\text{cl paths}}
D^{1/2}
e^{iS/\hbar - i\mu\pi/2}.
\]

%%%%%%%%%%%%%%%%%%%%%%%%%%%%%%%%%%%%%%%%%%%%%%%%%%%%%%%%%%%%
\chapter{Gutzwiller Trace Formula}

\section{Density of States}

\[
\rho(E)=\sum_n \delta(E-E_n).
\]

\section{Fourier Representation}

\[
\rho(E)
=
\frac1{2\pi\hbar}
\int e^{iEt/\hbar}
\text{Tr }U(t)dt.
\]

\section{Periodic Orbit Contribution}

\begin{theorem}[Gutzwiller Trace Formula]
\[
\rho(E)
=
\bar\rho(E)
+
\frac1{\pi\hbar}
\sum_\gamma
\frac{T_\gamma}
{\sqrt{|\det(M_\gamma-I)|}}
\cos\left(
\frac{S_\gamma(E)}{\hbar}
-
\frac{\pi\mu_\gamma}{2}
\right).
\]
\end{theorem}

\begin{proof}
Insert semiclassical propagator,
evaluate trace by stationary phase,
select periodic orbits,
Fourier transform in time.
\end{proof}

%%%%%%%%%%%%%%%%%%%%%%%%%%%%%%%%%%%%%%%%%%%%%%%%%%%%%%%%%%%%
\chapter{Quantum Chaos}

Level spacing distribution transitions:

\begin{itemize}
\item Poisson (integrable regime)
\item Wigner--Dyson (chaotic regime)
\end{itemize}

Periodic orbit proliferation explains spectral rigidity.

%%%%%%%%%%%%%%%%%%%%%%%%%%%%%%%%%%%%%%%%%%%%%%%%%%%%%%%%%%%%
\chapter{Outlook}

The double pendulum provides a unified framework linking:

\[
\text{Riemannian Geometry}
\rightarrow
\text{Hamiltonian Chaos}
\rightarrow
\text{Spectral Theory}
\rightarrow
\text{Quantum Chaos}.
\]

Future directions:

\begin{itemize}
\item Floquet driving
\item Microlocal analysis
\item Numerical periodic orbit theory
\item Experimental realizations
\end{itemize}

%%%%%%%%%%%%%%%%%%%%%%%%%%%%%%%%%%%%%%%%%%%%%%%%%%%%%%%%%%%%
\backmatter

\chapter*{Bibliographical Notes}

Classical mechanics: Arnold.  
Symplectic integration: Hairer, Lubich, Wanner.  
Quantum chaos: Gutzwiller.  

\end{document}
